% ============================================
% 表格模板 (Table Templates)
% 使用三线表规范，需 booktabs 宏包
% ============================================

% --------------------------------------------
% 基础模板
% --------------------------------------------
\begin{table}[H]
  \centering
  \caption{表格标题}
  \label{tab:xxx}
  \begin{tabular}{lcc}
    \toprule
    \textbf{列1} & \textbf{列2} & \textbf{列3} \\
    \midrule
    内容 & 内容 & 内容 \\
    内容 & 内容 & 内容 \\
    \bottomrule
  \end{tabular}
\end{table}

% --------------------------------------------
% 列对齐说明
% --------------------------------------------
% l - 左对齐（文本）
% c - 居中（数值、短文本）
% r - 右对齐（数字）
% p{宽度} - 固定宽度，自动换行

% --------------------------------------------
% 两列表格
% --------------------------------------------
\begin{table}[H]
  \centering
  \caption{缩略语对照表}
  \label{tab:terminology}
  \begin{tabular}{ll}
    \toprule
    \textbf{缩略语} & \textbf{全称} \\
    \midrule
    GDP  & 国内生产总值 \\
    CPI  & 消费者价格指数 \\
    R\&D & 研究与开发 \\
    \bottomrule
  \end{tabular}
\end{table}

% --------------------------------------------
% 数值对比表格
% --------------------------------------------
\begin{table}[H]
  \centering
  \caption{各组样本基本情况}
  \label{tab:comparison}
  \begin{tabular}{lcc}
    \toprule
    \textbf{组别} & \textbf{样本量} & \textbf{平均年龄} \\
    \midrule
    实验组 & 120 & 32.5 \\
    对照组 & 115 & 31.8 \\
    空白组 & 118 & 33.1 \\
    \bottomrule
  \end{tabular}
\end{table}

% --------------------------------------------
% 包含代码的表格
% --------------------------------------------
\begin{table}[H]
  \centering
  \caption{常用命令说明}
  \label{tab:commands}
  \begin{tabular}{lll}
    \toprule
    \textbf{命令} & \textbf{示例} & \textbf{功能} \\
    \midrule
    \texttt{git clone}  & \texttt{git clone <url>}   & 克隆远程仓库 \\
    \texttt{git commit} & \texttt{git commit -m "msg"} & 提交更改 \\
    \texttt{git push}   & \texttt{git push origin main} & 推送到远程 \\
    \bottomrule
  \end{tabular}
\end{table}

% --------------------------------------------
% 多行文本表格（使用 p{} 列）
% --------------------------------------------
\begin{table}[H]
  \centering
  \caption{研究方法特点对比}
  \label{tab:methods}
  \begin{tabular}{lp{8cm}}
    \toprule
    \textbf{方法} & \textbf{特点} \\
    \midrule
    问卷调查 & 成本较低，可大规模收集数据，但受访者可能存在
               主观偏差，问卷设计需严谨 \\
    深度访谈 & 能获取丰富的质性资料，深入了解受访者观点，
               但样本量有限，耗时较长 \\
    实地观察 & 可获得第一手真实数据，减少被试效应，
               但观察者可能影响被观察对象行为 \\
    \bottomrule
  \end{tabular}
\end{table}

% --------------------------------------------
% 带分组的表格（使用 \cmidrule）
% --------------------------------------------
\begin{table}[H]
  \centering
  \caption{评价指标体系}
  \label{tab:indicators}
  \begin{tabular}{llcc}
    \toprule
    \textbf{维度} & \textbf{指标} & \textbf{权重} & \textbf{数据来源} \\
    \midrule
    经济效益 & 投资回报率   & 0.25 & 财务报表 \\
    经济效益 & 成本节约率   & 0.20 & 财务报表 \\
    \cmidrule{1-4}
    社会效益 & 满意度评分   & 0.30 & 问卷调查 \\
    社会效益 & 覆盖人数     & 0.25 & 统计数据 \\
    \bottomrule
  \end{tabular}
\end{table}